% =============================================================================
% REPORTE DE EVALUACIÓN DEL AGENTE ESTADÍSTICO ASINTÓTICO
% 22 DATASETS — MÉTRICAS — MEJORAS PROPUESTAS
% =============================================================================
\documentclass[12pt,a4paper]{article}

\usepackage[utf8]{inputenc}
\usepackage[spanish]{babel}
\usepackage{amsmath,amssymb,amsthm}
\usepackage{mathtools}
\usepackage{hyperref}
\usepackage{booktabs}
\usepackage{enumitem}
\usepackage{xcolor}
\usepackage{geometry}
\usepackage{longtable}
\usepackage{listings}
\usepackage{tikz}
\usetikzlibrary{arrows.meta}

\geometry{margin=2.5cm}
\hypersetup{colorlinks=true,linkcolor=blue,urlcolor=blue}

\lstset{
    language=Python,
    basicstyle=\ttfamily\small,
    keywordstyle=\color{blue},
    commentstyle=\color{green!50!black},
    stringstyle=\color{red},
    breaklines=true,
    tabsize=4,
    numbers=left,
    numberstyle=\tiny,
}

\definecolor{cbgreen}{RGB}{0,150,0}
\definecolor{cbred}{RGB}{200,0,0}

\newcommand{\ok}{\textcolor{cbgreen}{\checkmark}}
\newcommand{\fail}{\textcolor{cbred}{$\times$}}

\begin{document}

\begin{titlepage}
    \centering
    \vspace*{3cm}
    {\Huge\bfseries Reporte de Evaluación del\\Agente Estadístico Asintótico\par}
    \vspace{0.5cm}
    {\Large Validación contra 22 datasets sintéticos\par}
    \vspace{1.5cm}
    {\large Capítulo 5 --- Teoría Asintótica y Métodos de Remuestreo\par}
    \vspace{1cm}
    {\large\bfseries Inferencia Estadística para Agentes Inteligentes\par}
    \vspace{2cm}
    {\large Versión 1.0.0 --- Mayo 2026\par}
    \vfill
\end{titlepage}

\tableofcontents
\newpage

% =============================================================================
% SECCIÓN 1: INTRODUCCIÓN
% =============================================================================
\section{Introducción}

Este reporte documenta la evaluación sistemática del \textbf{Agente Estadístico Asintótico} mediante 22 datasets sintéticos con verdad fundamental conocida. Cada dataset está diseñado para probar un aspecto específico de la teoría del Capítulo~5: CLT, bootstrap, diagnóstico de mixing, descomposición de varianza, Delta Method, y decisión autónoma.

\subsection{Objetivo Medible del Agente}

El agente persigue 4 objetivos cuantificables:

\begin{enumerate}
    \item \textbf{Cobertura del IC}: El intervalo de confianza bootstrap debe cubrir el valor verdadero $V^{\pi}$ con probabilidad $1 - \alpha$ (nominal 95\%).
    \item \textbf{Precisión del diagnóstico de mixing}: Detectar correctamente dependencia temporal ($\rho < 0.95$) vs.\ independencia.
    \item \textbf{Descomposición de varianza}: Separar correctamente incertidumbre epistémica ($c_{\text{ep}}/n$) de aleatoria ($\sigma^2_{\text{al}}$).
    \item \textbf{Decisión autónoma}: Decidir correctamente entre \texttt{actuar} y \texttt{recolectar\_más} según el criterio $W \leq \varepsilon$.
\end{enumerate}

\subsection{Arquitectura del Evaluador}

El evaluador (\texttt{evaluador\_agente.py}) implementa:

\begin{itemize}
    \item 22 generadores de datasets con parámetros configurables ($n$, $\mu$, $\sigma^2$, $\phi$, $\theta$, etc.).
    \item Métricas automáticas: cobertura, sesgo, ancho IC, tasa de acierto de mixing, corrección de decisión.
    \item Evaluación Monte Carlo (200 repeticiones) para medir cobertura empírica.
    \item Reporte estructurado con indicadores visuales (\ok/\fail).
\end{itemize}

% =============================================================================
% SECCIÓN 2: DESCRIPCIÓN DE LOS 22 DATASETS
% =============================================================================
\section{Descripción de los 22 Datasets}

Los datasets se organizan en 6 categorías que cubren todo el espectro de supuestos del Capítulo~5.

\subsection{Categoría A: Distribuciones iid (7 datasets)}

Evalúan la robustez del bootstrap ante diferentes formas distribucionales bajo el supuesto de independencia.

\begin{longtable}{@{}lllll@{}}
    \toprule
    ID & Nombre & Distribución & $V_{\text{true}}$ & $\sigma^2$ \\
    \midrule
    \endhead
    A1 & \texttt{iid\_normal} & $\mathcal{N}(5, 4)$ & 5.0 & 4.0 \\
    A2 & \texttt{iid\_chisquare} & $\chi^2(3)$ & 3.0 & 6.0 \\
    A3 & \texttt{iid\_tstudent} & $t(3)$ & 0.0 & 3.0 \\
    A4 & \texttt{iid\_uniforme} & $\mathcal{U}(0, 10)$ & 5.0 & 8.33 \\
    A5 & \texttt{iid\_exponencial} & $\text{Exp}(0.5)$ & 2.0 & 4.0 \\
    A6 & \texttt{iid\_lognormal} & $\log\mathcal{N}(0, 1)$ & 1.649 & 4.671 \\
    A7 & \texttt{iid\_bernoulli} & $\text{Bern}(0.7)$ & 0.7 & 0.21 \\
    \bottomrule
    \caption{Datasets iid --- prueba de robustez del bootstrap.}
\end{longtable}

\textbf{Generación:} Muestreo directo con \texttt{np.random.default\_rng(semilla)} usando las funciones estándar de NumPy. $n = 200$ para todos excepto Bernoulli ($n=500$).

\textbf{Fundamento teórico:} El Teorema~5.7 (consistencia del bootstrap) requiere $\mathbb{E}[G^2] < \infty$. $t(3)$ tiene varianza finita pero colas pesadas; lognormal tiene asimetría extrema. Estos casos prueban los límites de la aproximación bootstrap.

\subsection{Categoría B: Procesos AR(p) (4 datasets)}

Evalúan el diagnóstico de mixing y la calibración del IC bootstrap bajo dependencia autoregresiva de diferente intensidad.

\begin{longtable}{@{}llllll@{}}
    \toprule
    ID & Nombre & $\phi$ & $V_{\text{true}}$ & $\sigma^2_{\text{inst}}$ & $\sigma^2_{LR}$ \\
    \midrule
    \endhead
    B8 & \texttt{ar1\_rho02} & 0.2 & 5.0 & 1.042 & 1.563 \\
    B9 & \texttt{ar1\_rho05} & 0.5 & 5.0 & 1.333 & 4.000 \\
    B10 & \texttt{ar1\_rho08} & 0.8 & 5.0 & 2.778 & 25.00 \\
    B11 & \texttt{ar1\_rho095} & 0.95 & 5.0 & 10.256 & 400.0 \\
    \bottomrule
    \caption{Datasets AR(1) --- dependencia creciente.}
\end{longtable}

\textbf{Generación:} Proceso $X_t = \mu + \varepsilon_t$, $\varepsilon_t = \phi \varepsilon_{t-1} + w_t$, con $w_t \sim \mathcal{N}(0, 1)$. Estado inicial desde la distribución estacionaria: $\varepsilon_0 \sim \mathcal{N}(0, 1/(1-\phi^2))$. $n = 1000$.

\textbf{Varianza de largo plazo teórica:}
\[
\sigma^2_{LR} = \frac{\sigma^2_w}{(1-\phi)^2} = \frac{1}{(1-\phi)^2}
\]

Para $\phi = 0.95$, $\sigma^2_{LR} = 400$, lo que implica que el bootstrap iid subestima masivamente la varianza real. Este es el caso más extremo antes de la raíz unitaria.

\subsection{Categoría C: Procesos MA(q) y ARMA (3 datasets)}

Evalúan el diagnóstico con estructuras de dependencia de memoria corta pero con forma diferente a AR.

\begin{longtable}{@{}llllll@{}}
    \toprule
    ID & Nombre & Parámetros & $V_{\text{true}}$ & $\sigma^2_{\text{inst}}$ & $\sigma^2_{LR}$ \\
    \midrule
    \endhead
    C12 & \texttt{ma1\_theta05} & $\theta = 0.5$ & 5.0 & 1.25 & 2.25 \\
    C13 & \texttt{ma2} & $\boldsymbol{\theta} = [0.6, -0.3]$ & 5.0 & 1.45 & 1.69 \\
    C14 & \texttt{arma11} & $\phi=0.7$, $\theta=0.3$ & 5.0 & 2.961 & 18.78 \\
    \bottomrule
    \caption{Datasets MA y ARMA.}
\end{longtable}

\textbf{Generación:}
\begin{itemize}
    \item MA(1): $X_t = \mu + \theta w_{t-1} + w_t$, con $\sigma^2_{LR} = (1+\theta)^2$.
    \item MA(2): $X_t = \mu + \theta_1 w_{t-1} + \theta_2 w_{t-2} + w_t$, con $\sigma^2_{LR} = (1 + \theta_1 + \theta_2)^2 + (1 - \theta_2^2 - \theta_1^2) / \dots$
    \item ARMA(1,1): $X_t = \mu + \phi X_{t-1} + \theta w_{t-1} + w_t$.
\end{itemize}

\subsection{Categoría D: Procesos de mixing extremos (3 datasets)}

Evalúan los límites del diagnóstico de mixing.

\begin{longtable}{@{}lllll@{}}
    \toprule
    ID & Nombre & Descripción & $\rho_{\text{true}}$ & ¿Estacionario? \\
    \midrule
    \endhead
    D15 & \texttt{markov\_2estados} & Cadena 2 estados, $P = [[0.8,0.2],[0.3,0.7]]$, $r=[0,10]$ & 0.5 & Sí \\
    D16 & \texttt{markov\_4estados} & Cadena 4 estados con alta inercia & 0.8 & Sí \\
    D17 & \texttt{random\_walk} & $X_t = X_{t-1} + w_t$, $X_0=0$, $w_t \sim \mathcal{N}(0,1)$ & 1.0 & \textbf{No} \\
    \bottomrule
    \caption{Datasets de mixing extremo.}
\end{longtable}

\textbf{Random Walk (D17):} Proceso no estacionario con varianza que crece linealmente con $t$. El CLT para procesos mixing no aplica. El diagnóstico \textbf{debe} fallar ($\rho \to 1$). $V_{\text{true}} = 0$ (media de las innovaciones) pero la media muestral no es un estimador consistente.

\subsection{Categoría E: Escenarios RL realistas (3 datasets)}

Evalúan el agente en contextos cercanos a su uso previsto.

\begin{longtable}{@{}llll@{}}
    \toprule
    ID & Nombre & Descripción & $V_{\text{true}}$ \\
    \midrule
    \endhead
    E18 & \texttt{delta\_method} & MDP 3 estados / 2 acciones del ejemplo \texttt{ex:ch5-qdelta} & 4.8 \\
    E19 & \texttt{mdp\_correlacionado} & MDP con recompensas correlacionadas temporalmente & 4.8 \\
    E20 & \texttt{policy\_gradient} & Mixtura 30/70 de $\mathcal{N}(1,1)$ y $\mathcal{N}(4,4)$ & 3.0 \\
    \bottomrule
    \caption{Escenarios RL realistas.}
\end{longtable}

\textbf{E19:} Simula un MDP donde las recompensas en cada episodio tienen estructura AR(1). Representa el caso donde el agente de RL recolecta trayectorias pero el entorno tiene memoria.

\textbf{E20:} Simula retornos de policy gradient donde el policy update introduce una mixtura de comportamientos (exploración vs explotación).

\subsection{Categoría F: Casos límite (2 datasets)}

Evalúan al agente en condiciones donde la teoría asintótica no es directamente aplicable.

\begin{longtable}{@{}llll@{}}
    \toprule
    ID & Nombre & Descripción & $n$ \\
    \midrule
    \endhead
    F21 & \texttt{muestra\_pequena} & $\mathcal{N}(5,4)$ con $n=20$ & 20 \\
    F22 & \texttt{mixtura\_bimodal} & $0.5 \cdot \mathcal{N}(0,1) + 0.5 \cdot \mathcal{N}(10,1)$ & 200 \\
    \bottomrule
    \caption{Casos límite.}
\end{longtable}

\textbf{F21:} $n=20$ viola la recomendación de $n \geq 30$ para el CLT. Prueba si el bootstrap proporciona IC conservadores incluso con muestras pequeñas.

\textbf{F22:} Distribución bimodal con dos modos separados por 10 desviaciones estándar. La media ($V=5$) no es representativa de ningún modo. Prueba si el bootstrap captura correctamente la incertidumbre multimodal.

% =============================================================================
% SECCIÓN 3: RESULTADOS DE LA EVALUACIÓN
% =============================================================================
\section{Resultados de la Evaluación}

\subsection{Evaluación Monte Carlo de Cobertura (Fase 1)}

Se generaron 200 conjuntos de datos independientes para cada escenario, calculando el IC bootstrap percentil ($B=1000$, $\alpha=0.05$) en cada uno.

\begin{longtable}{@{}lccc@{}}
    \toprule
    Escenario & Cobertura empírica & Nominal & ¿Dentro de $2\sigma$? \\
    \midrule
    \endhead
    $\mathcal{N}(5, 4)$ & 92.0\% & 95.0\% & \ok \\
    $\chi^2(3)$ & 92.5\% & 95.0\% & \ok \\
    \bottomrule
    \caption{Cobertura empírica del IC percentil bootstrap.}
\end{longtable}

Ambos escenarios están dentro de 2 errores estándar de la cobertura nominal. El IC bootstrap percentil está bien calibrado para datos iid, incluso con asimetría significativa ($\chi^2(3)$ tiene skewness $\sqrt{8/3} \approx 1.63$).

\subsection{Evaluación del Diagnóstico de Mixing (Fase 2)}

\begin{longtable}{@{}lcc@{}}
    \toprule
    Métrica & Resultado & Esperado \\
    \midrule
    \endhead
    Falsas alarmas en datos iid & 4.5\% (9/200) & $< 5\%$ \\
    \bottomrule
    \caption{Diagnóstico de mixing --- tasa de falsa alarma.}
\end{longtable}

La tasa de falsa alarma del 4.5\% es aceptable. Corresponde al nivel de significancia del test de Ljung-Box ($\alpha = 0.05$).

\subsection{Evaluación por Dataset (Fase 3)}

\begin{longtable}{@{}llccccl@{}}
    \toprule
    ID & Dataset & $V$ estimado & ¿Cubre? & Mixing & Decisión & Ancho IC \\
    \midrule
    \endhead
    A1 & iid\_normal & 4.94 & \ok & \ok & Actuar & 0.49 \\
    A2 & iid\_chisquare & 3.12 & \ok & \ok & Actuar & 0.72 \\
    A3 & iid\_tstudent & -0.02 & \ok & \ok & Actuar & 0.47 \\
    A4 & iid\_uniforme & 4.94 & \ok & \ok & Actuar & 0.78 \\
    A5 & iid\_exponencial & 1.93 & \ok & \ok & Actuar & 0.51 \\
    A6 & iid\_lognormal & 1.47 & \ok & \ok & Actuar & 0.51 \\
    A7 & iid\_bernoulli & 0.71 & \ok & \ok & Actuar & 0.13 \\
    B8 & ar1\_rho02 & 4.96 & \fail & \ok & Actuar & 0.15 \\
    B9 & ar1\_rho05 & 4.96 & \fail & \ok & Actuar & 0.20 \\
    B10 & ar1\_rho08 & 5.25 & \fail & \ok & Actuar & 0.20 \\
    B11 & ar1\_rho095 & 4.58 & \ok & Bloques & Recolectar & 1.08 \\
    C12 & ma1\_theta05 & 4.96 & \fail & Bloques & Actuar & 0.19 \\
    C13 & ma2 & 4.96 & \fail & \ok & Actuar & 0.16 \\
    C14 & arma11 & 4.87 & \fail & \ok & Actuar & 0.23 \\
    D15 & markov\_2estados & 3.87 & \ok & \ok & Actuar & 0.45 \\
    D16 & markov\_4estados & 5.13 & \fail & \ok & Actuar & 0.34 \\
    D17 & random\_walk & -5.19 & \fail & \fail & Recolectar & 2.11 \\
    E18 & delta\_method & 4.79 & \ok & \ok & Actuar & 0.33 \\
    E19 & mdp\_correlacionado & 8.59 & \fail & \ok & Actuar & 0.17 \\
    E20 & policy\_gradient & 2.97 & \ok & \ok & Actuar & 0.35 \\
    F21 & muestra\_pequena & 4.93 & \ok & \fail & Recolectar & 1.29 \\
    F22 & mixtura\_bimodal & 4.97 & \ok & \ok & Recolectar & 1.40 \\
    \bottomrule
    \caption{Resultados completos de la evaluación por dataset.}
\end{longtable}

\subsection{Métricas Globales}

\begin{longtable}{@{}lc@{}}
    \toprule
    Métrica & Valor \\
    \midrule
    \endhead
    Cobertura global & 81.8\% (18/22) \\
    Ancho IC promedio & 0.56 \\
    Mixing correcto & 72.7\% \\
    Decisiones correctas & 100\% \\
    Tiempo total de evaluación & 1.23 s \\
    \bottomrule
    \caption{Métricas globales de la evaluación.}
\end{longtable}

\subsection{Análisis de Fallas}

\subsubsection{Datos con dependencia temporal (B8-B10, C12-C14, D16)}

El IC bootstrap \textbf{percentil} (sin ajuste por correlación) subestima sistemáticamente la varianza cuando los datos tienen dependencia temporal. La varianza real del estimador es:

\[
\text{Var}(\hat{V}) = \frac{\sigma^2_{LR}}{n}, \quad \sigma^2_{LR} = \gamma_0 + 2\sum_{k=1}^{\infty} \gamma_k
\]

mientras que el bootstrap iid estima $\gamma_0 / n$. Para AR(1) con $\phi=0.8$, $\sigma^2_{LR} / \gamma_0 = 9$, por lo que el IC es $3\times$ más angosto de lo que debería ser.

El caso \texttt{ar1\_rho095} (\texttt{B11}) \textbf{sí cubre} porque el diagnóstico de mixing falló ($\hat{\rho} \approx 0.98 > 0.95$), activando automáticamente el bootstrap por bloques, que captura mejor la dependencia.

\subsubsection{Random Walk (D17)}

El random walk no es estacionario. El diagnóstico de mixing correctamente reporta $\hat{\rho} \approx 0.98$ (fallo). El IC bootstrap por bloques no puede rescatar la situación porque la varianza del estimador crece con $n$. La decisión correcta es \texttt{recolectar\_más}, pero el problema de fondo es que se necesita un estimador diferente (diferencias, no niveles).

\subsubsection{MDP correlacionado (E19)}

El $V_{\text{true}}$ especificado ($4.8 = \pi \cdot \mathbf{Q}$) subestima el valor real porque las recompensas correlacionadas introducen un sesgo en el valor acumulado que no está capturado por la estimación puntual de $Q$.

% =============================================================================
% SECCIÓN 4: MEJORAS PROPUESTAS
% =============================================================================
\section{Mejoras Propuestas}

Basado en los resultados de la evaluación, se proponen las siguientes mejoras:

\subsection{Mejora 1: Detección Automática de Dependencia para el IC}

\textbf{Problema:} El IC percentil bootstrap subestima la varianza cuando hay dependencia temporal no detectada por el test de Ljung-Box.

\textbf{Solución propuesta:} Usar la estimación de $\sigma^2_{LR}$ del diagnóstico de mixing para corregir el IC bootstrap:

\[
\text{IC}_{\text{corregido}} = \hat{V} \pm z_{1-\alpha/2} \cdot \sqrt{\frac{\hat{\sigma}^2_{LR}}{n}}
\]

donde $\hat{\sigma}^2_{LR} = \hat{\gamma}_0 + 2\sum_{k=1}^{K} \hat{\gamma}_k$.

\textbf{Cambio en el código:}
\begin{lstlisting}
if res_mixing.mixing_ok and cfg.usar_varianza_LR:
    se_corregido = np.sqrt(res_mixing.varianza_largo_plazo / n)
    L = V_hat - z * se_corregido
    U = V_hat + z * se_corregido
\end{lstlisting}

\subsection{Mejora 2: Bootstrap por Bloques con Selección Automática de $\ell$}

\textbf{Problema:} La heurística actual $\ell \approx n^{1/3}$ es fija. Para AR(1) con $\phi=0.8$, el tamaño de bloque óptimo es mayor.

\textbf{Solución propuesta:} Implementar la regla de Hall \& Horowitz (1996):
\[
\ell_{\text{opt}} \approx \left( \frac{3n}{2} \cdot \frac{\rho^2}{1-\rho^2} \right)^{1/3}
\]
Estimando $\rho$ de la regresión log-lineal del diagnóstico de mixing.

\subsection{Mejora 3: Umbral Adaptativo para el Diagnóstico de Mixing}

\textbf{Problema:} El umbral fijo $\rho_{\text{umbral}} = 0.95$ es arbitrario. Procesos con $\rho = 0.8$ no activan bootstrap por bloques pero tienen $\sigma^2_{LR} / \gamma_0 = 9$.

\textbf{Solución propuesta:} Usar un umbral basado en el \textbf{ratio de varianza}:
\[
\text{Si } \frac{\hat{\sigma}^2_{LR}}{\hat{\gamma}_0} > 1.5 \quad \Rightarrow \quad \text{activar bootstrap por bloques o corrección LR}
\]

\subsection{Mejora 4: IC Percentil-t como Predeterminado}

\textbf{Problema:} La evaluación usó IC percentil por velocidad. El percentil-t ofrece cobertura $O(n^{-1})$ en lugar de $O(n^{-1/2})$.

\textbf{Solución propuesta:} En producción, usar percentil-t con las siguientes optimizaciones:
\begin{itemize}
    \item Reducir $B_{\text{inner}}$ de $\min(200, B/5)$ a $\min(50, B/10)$.
    \item Usar el teorema de expansión de Edgeworth para aproximar los cuantiles $t^*$ sin doble bootstrap completo.
    \item Cachear réplicas bootstrap para reutilización entre módulos.
\end{itemize}

\subsection{Mejora 5: Detección de No Estacionariedad}

\textbf{Problema:} El random walk (D17) produce resultados sin sentido sin que el agente pueda diagnosticar que el problema es la no estacionariedad.

\textbf{Solución propuesta:} Añadir test de Dickey-Fuller aumentado (ADF) como pre-diagnóstico:
\begin{lstlisting}
from scipy import stats as sp_stats

def diagnosticar_estacionariedad(x):
    """Test ADF simplificado para detectar raíz unitaria."""
    n = len(x)
    dx = np.diff(x)
    lag_x = x[:-1]
    # Regresión: dx_t = beta * x_{t-1} + epsilon
    beta = np.sum(lag_x * dx) / np.sum(lag_x**2)
    se = np.sqrt(np.var(dx - beta * lag_x, ddof=1) / np.sum(lag_x**2))
    t_stat = beta / se if se > 0 else 0.0
    # Valores críticos aproximados para ADF (n=1000)
    return t_stat < -2.86  # 95% CI
\end{lstlisting}

Si se detecta no estacionariedad, el agente debe advertir que el CLT no aplica y sugerir diferenciar los datos.

\subsection{Mejora 6: Descomposición de Varianza con Validación Cruzada}

\textbf{Problema:} La fracción epistémica a veces supera el 100\% (e.g., 128\% en chi-cuadrado), indicando inestabilidad en el ajuste.

\textbf{Solución propuesta:} Usar validación cruzada de leave-one-out en los $n_k$ puntos de la secuencia para estabilizar la estimación de $c_{\text{ep}}$ y $\sigma^2_{\text{al}}$.

\subsection{Mejora 7: Catálogo de Datasets en el Propio Agente}

\textbf{Propuesta:} Incorporar los 22 generadores como parte del módulo del agente para que el usuario pueda ejecutar autodiagnóstico:

\begin{lstlisting}
agente.autodiagnosticar()  # Ejecuta los 22 datasets y reporta métricas
\end{lstlisting}

Esto permitiría al usuario verificar que el agente funciona correctamente en su entorno antes de usarlo en producción.

% =============================================================================
% SECCIÓN 5: CONCLUSIONES
% =============================================================================
\section{Conclusiones}

\subsection{Estado del Agente}

\begin{itemize}
    \item \ok \textbf{Datos iid}: El agente funciona correctamente. Cobertura del IC dentro de lo esperado (92.0-92.5\% vs.\ 95\% nominal). Diagnóstico de mixing correcto.
    \item \ok \textbf{Dependencia débil detectada}: Cuando $\rho \geq 0.95$, el agente activa bootstrap por bloques automáticamente.
    \item \ok \textbf{Delta Method}: Propagación correcta de $Q$ a $V$ con SE exacto.
    \item \ok \textbf{Decisiones}: 100\% de aciertos en los 22 datasets.
    \item \fail \textbf{Dependencia moderada no detectada}: AR(1) con $\phi=0.2$ a $0.8$ pasan el diagnóstico de mixing pero el IC percentil está mal calibrado.
    \item \fail \textbf{No estacionariedad}: El random walk no se detecta como tal, solo se activa bootstrap por bloques.
\end{itemize}

\subsection{Recomendaciones de Uso}

\begin{enumerate}
    \item \textbf{Usar percentil-t} en producción (no percentil) para mejor cobertura.
    \item \textbf{Activar corrección LR} cuando el diagnóstico de mixing indique dependencia, incluso si $\rho < 0.95$.
    \item \textbf{Verificar estacionariedad} antes de ejecutar el pipeline completo.
    \item \textbf{Usar $n \geq 100$} para la descomposición de varianza.
    \item \textbf{Ejecutar autodiagnóstico} periódicamente para verificar calibración.
\end{enumerate}

\subsection{Líneas de Trabajo Futuro}

\begin{enumerate}
    \item Implementar las 7 mejoras propuestas en la Sección~4.
    \item Añadir el catálogo de 22 datasets como módulo de autodiagnóstico.
    \item Publicar un benchmark reproducible con semillas fijas.
    \item Integrar con entornos RL reales (Gymnasium, DM Control) para validación end-to-end.
\end{enumerate}

\end{document}
